% Insert after the pilot setup and metric definitions.
% Evidence: normalized_improvement_pilots_v1; no new experiments.
% Include the two normalized-improvement table files for cross-references.
\subsection{Pilot Findings and Motivation}
\label{sec:pilot-findings}

\paragraph{Reporting convention.}
Both pilots use the same normalized improvement score:
\begin{equation}
  \bar I = \frac{1}{N}\sum_{i=1}^{N}\frac{1}{3}\sum_{s=0}^{2}
  \frac{10\bigl(S^{\mathrm{opt}}_{i,s}-S^{\mathrm{orig}}_{i,s}\bigr)}
       {S^{\mathrm{orig}}_{i,s}+1}.
  \label{eq:pilot-normalized-improvement}
\end{equation}
For each model and strategy, we first compute improvement for matched
instance--answer-seed pairs, average over three seeds within each instance,
and then average over the $N$ available instances in each benchmark.
Objective uses zero-fallback GEO Overall (base scores 0--1), and Subjective
uses the RAID seven-dimension average (base scores 0--5).
Normalized improvement scores are not percentages and need not lie within
these base-score ranges. Zero Original scores are retained under the same
formula. Boldface denotes the highest mean in the stated candidate set,
not statistical significance.

\paragraph{Prompt-family pilot.}
Table~\ref{tab:necessity-part1-normalized-improvement}
shows that the relative ordering of some optimization prompts varies across
answer models. On GEO-Bench, Conclusion-First achieves a larger objective
normalized improvement score than Hierarchical Framing on Qwen
(0.0794 versus 0.0191), whereas Hierarchical Framing scores higher on
Llama (0.1131 versus 0.0800) and Mistral (0.3397 versus 0.1647).
The corresponding subjective scores exhibit the same ordering reversal:
2.5574 versus 1.7545 on Qwen, 2.2818 versus 1.9225 in favor of
Hierarchical Framing on Llama, and 3.6535 versus 2.3297 on Mistral.
These local differences coexist with a shared highest-scoring strategy.
Key-Point Enumeration achieves the highest mean among the nine optimization
prompts for both metrics, all three models, and both benchmarks.

\paragraph{Literature-method pilot.}
Table~\ref{tab:necessity-part2-normalized-improvement}
shows analogous differences for literature-derived methods under our
experimental configuration. On C-SEO Bench, Statistics Addition exceeds
Content Improvement in objective normalized improvement on Qwen
(0.1228 versus 0.0230). The ordering reverses on Llama
(0.2663 versus 0.0577 in favor of Content Improvement) and Mistral
(0.2610 versus 0.0987). Nevertheless, LLM Guidance attains the highest
mean among the five fixed methods for every model, benchmark, and metric.
We report RAID-GSEO and IF-GEO separately from this fixed-method comparison
because their multi-step procedures have different computational budgets.
IF-GEO additionally has incomplete sample coverage, so its reported means
describe successful rewrites and their matched Original conditions.

\paragraph{Implications for model-specific adaptation.}
Together, the two pilots provide descriptive evidence that the relative
effectiveness of a given GEO strategy can depend on the downstream answer
model. This observation motivates studying whether a given optimization
prompt can be adapted to improve its effectiveness on a target model.
The pilots do not establish that adaptation is necessary or that a
model-specific strategy outperforms a strong fixed alternative.
Indeed, the shared highest-scoring fixed strategies in both candidate sets
provide important counterpoints to that stronger claim. Demonstrating the
value of adaptation requires comparisons against direct application of
the original strategy and strong fixed-strategy baselines.

\paragraph{Scope of interpretation.}
These are descriptive comparisons of mean normalized improvement scores,
not tests of statistical significance or transfer performance. Both pilots
use the same paired-seed computation and aggregation procedure. Their
available sample sets differ, and source-prefix truncation affects
16 conditions in the literature-method pilot. The shared GLM rewriter
checkpoint does not imply identical generation implementations across
the two runs, including their end-of-sequence handling. These differences
limit direct cross-pilot comparisons and attribution solely to rewriting.
Furthermore, normalized improvement weights changes by the Original score:
a positive mean does not necessarily imply an increase in mean absolute
visibility, and stochastic variation can contribute to positive values.
We therefore retain the Original and Neutral Rewrite controls where
available, together with the underlying visibility scores, when interpreting
the observed gains.
